\documentclass{article}

\usepackage[utf8]{inputenc}
\usepackage{geometry}
\usepackage{listings}
\usepackage{xcolor}

\geometry{margin=1in}

\title{Reporte de como se usa matplotlib y pandas}
\author{Reporte de Actividad}
\date{\today}

\lstset{
    language=Python,
    basicstyle=\ttfamily\small,
    keywordstyle=\color{blue},
    commentstyle=\color{gray},
    stringstyle=\color{green!50!black},
    breaklines=true
}

\begin{document}

\maketitle

\section{Introducción}

En esta actividad se realizó un análisis exploratorio de datos utilizando las librerías
\textbf{Pandas} y \textbf{Matplotlib} en Python, el dataset utilizado contiene información
sobre viviendas, incluyendo variables como precio, número de habitaciones,
pisos y otras características estructurales.

El objetivo fue explorar los datos com las librerías de python para poder , calcular estadísticas básicas y generar
visualizaciones para identificar patrones en la información imprencidible en ia.

\section{Carga de librerías}

Primero se importaron las librerías necesarias para el análisis.

\begin{lstlisting}
import pandas as pd
import matplotlib.pyplot as plt
\end{lstlisting}

\section{Carga del dataset}

Se utilizó la función \texttt{read\_csv()} de Pandas para cargar el archivo
\texttt{housedata\_audit\_v1.csv} en un DataFrame.

\begin{lstlisting}
df = pd.read_csv("housedata_audit_v1.csv")
df.head()
\end{lstlisting}

Esta instrucción permite visualizar las primeras filas del dataset.

\section{Exploración de los datos}

Para conocer la estructura del dataset se utilizaron varias funciones de Pandas.

\begin{lstlisting}
df.shape
df.columns
df.info()
\end{lstlisting}

\begin{itemize}
\item \textbf{df.shape} muestra el número de filas y columnas.
\item \textbf{df.columns} muestra los nombres de las variables.
\item \textbf{df.info()} proporciona información sobre los tipos de datos.
\end{itemize}

\section{Estadísticas descriptivas}

Se calcularon estadísticas básicas de las variables numéricas utilizando:

\begin{lstlisting}
df.describe()
\end{lstlisting}

Esto permite obtener información como:

\begin{itemize}
\item media
\item desviación estándar
\item valor mínimo
\item valor máximo
\end{itemize}

\section{Análisis de valores nulos}

Para identificar datos faltantes se utilizó:

\begin{lstlisting}
df.isnull().sum()
\end{lstlisting}

Esto muestra cuántos valores faltantes existen en cada columna.

\section{Visualización de datos}

Se generaron varias gráficas utilizando Matplotlib.

\subsection{Histograma del precio}

\begin{lstlisting}
plt.figure()

df['Meta_Dato_01'].hist(bins=30)

plt.title("Distribución del precio de casas")
plt.xlabel("Precio")
plt.ylabel("Frecuencia")

plt.show()
\end{lstlisting}

Este gráfico permite observar la distribución de los precios de las viviendas.

\subsection{Relación entre habitaciones y precio}

\begin{lstlisting}
plt.figure()

plt.scatter(df['bedrooms'], df['Meta_Dato_01'])

plt.title("Habitaciones vs Precio")
plt.xlabel("Número de habitaciones")
plt.ylabel("Precio")

plt.show()
\end{lstlisting}

La gráfica de dispersión permite identificar si existe relación entre
el número de habitaciones y el precio de las casas.

\subsection{Precio promedio por número de pisos}

Primero se calculó el promedio de precio por número de pisos.

\begin{lstlisting}
precio_promedio = df.groupby('floors')['Meta_Dato_01'].mean()
\end{lstlisting}

Posteriormente se generó una gráfica de barras.

\begin{lstlisting}
plt.figure()

plt.bar(precio_promedio.index, precio_promedio.values)

plt.title("Precio promedio por número de pisos")
plt.xlabel("Pisos")
plt.ylabel("Precio promedio")

plt.show()
\end{lstlisting}

\section{Conclusiones}

El análisis exploratorio permitió identificar características importantes
del dataset de viviendas. Las herramientas de Pandas facilitaron la manipulación
y análisis de los datos, mientras que Matplotlib permitió visualizar patrones
como la distribución de precios y la relación entre distintas variables.

Estas técnicas son fundamentales en las etapas iniciales de proyectos de
aprendizaje automático y análisis de datos.

\end{document}
